\section{Related Work}
\label{sec:related-work}

\textbf{eBPF-based Monitoring.}
The rise of cloud-native computing has driven the adoption of eBPF for observability and security.
Pixie~\cite{pixie} uses eBPF for auto-telemetry, collecting rich application-level data with minimal overhead.
Falco~\cite{falco} and Tetragon~\cite{tetragon} leverage eBPF to monitor system calls and enforce security policies.
However, these systems primarily rely on manually defined rules or basic n-gram anomalies.
\textsc{WasmSentinel} extends this by enabling sandboxed ML inference directly on the eBPF event stream, bridging the gap between high-performance monitoring and advanced detection logic.

\textbf{Deep Learning for Intrusion Detection.}
Deep learning has shown great promise in detecting sophisticated attacks.
DeepLog~\cite{deeplog} utilizes LSTMs to model system logs as a natural language, detecting anomalies by predicting the next log entry.
Similarly, various architectures including Transformers~\cite{transformer-sec} and Graph Neural Networks (GNNs)~\cite{gnn-sec} have been applied to syscall sequences.
While accurate, these models often incur prohibitive inference overhead.
\textsc{WasmSentinel} demonstrates that a carefully optimized MLP, when combined with a hybrid architecture and Wasm-based execution, can achieve competitive accuracy with microsecond-level latency.

\textbf{WebAssembly for Security and Sandboxing.}
WebAssembly's strong isolation and near-native performance have made it an attractive target for security-sensitive extensions.
CT-Wasm~\cite{ct-wasm} proposes a type-driven approach to ensure constant-time execution in Wasm, mitigating side-channel attacks.
Other works have used Wasm to sandbox untrusted drivers or browser extensions.
In \textsc{WasmSentinel}, we leverage Wasm not only for isolation but as a portable, high-performance runtime for heterogeneous detection logic, enabling secure deployment of community-contributed ML models.
